% ======================================================
% 第 13 天：罗尔中值定理及其应用
% ======================================================
\setday{13}{罗尔中值定理及其应用}

\oneproblem{
    设函数 $f(x) = x^2(x-1)(x-2)$，则 $f'(x)$ 的零点个数为( )
    \begin{multicols}{4}
    \begin{enumerate}[label=(\Alph*)]
        \item 0 \item 1 \item 2 \item 3
    \end{enumerate}
    \end{multicols}
}

\oneproblem{
    设函数 $f(x)$ 在 $[a,b]$ 上可导，且 $f'_+(a) \ne f'_-(b)$，$\eta$ 是介于 $f'_+(a), f'_-(b)$ 之间的任一实数，则至少存在一点 $\xi \in (a,b)$，使得 $f'(\xi) = \eta$.
}

\oneproblem{
    设 $f(x)$ 在 $[0,1]$ 上可导，且 $f(1) = 2f(0)$. 求证存在 $\xi \in (0,1)$，使得 $(\xi + 1)f'(\xi) = f(\xi)$.
}

\oneproblem{
    设 $f(x)$ 在 $[0,1]$ 上二次可导，且 $f(0) = f(1) = 0$. 试证: 存在 $\xi \in (0,1)$，使得 $f''(\xi) = \dfrac{f'(\xi)}{(\xi - 1)^2}$.
}

\oneproblem{
    设 $f(x)$ 在 $[0,1]$ 上二阶可导，$f(0) = f(1)$. 求证: 至少存在一点 $\xi \in (0,1)$，使得 $2f'(\xi) + (\xi - 1)f''(\xi) = 0$.
}

\oneproblem{
    假设函数 $f(x)$ 和 $g(x)$ 在 $[a,b]$ 上存在二阶导数，并且 $g''(x) \ne 0$, $f(a) = f(b) = g(a) = g(b) = 0$. 试证:
    \begin{enumerate}[label=(\arabic*)]
        \item 在开区间 $(a,b)$ 内 $g(x) \ne 0$;
        \item 在开区间 $(a,b)$ 内至少存在一点 $\xi$，使得 $\dfrac{f(\xi)}{g(\xi)} = \dfrac{f''(\xi)}{g''(\xi)}$.
    \end{enumerate}
}

\oneproblem{
    设函数 $f(x)$ 在区间 $[a,b]$ 上具有二阶导数，且 $f(a) = f(b) = 0$, $f'_+(a)f'_-(b) > 0$. 证明存在 $\xi \in (a,b)$ 和 $\eta \in (a,b)$，使 $f(\xi) = 0$ 及 $f''(\eta) = 0$.
}

\oneproblem{
    设函数 $f(x)$ 在区间 $[0,1]$ 上连续，在 $(0,1)$ 内可导，$f(0) = f(1) = 0, f\left(\dfrac{1}{2}\right) = 1$，试证:
    \begin{enumerate}[label=(\arabic*)]
        \item 存在 $\eta \in \left(\dfrac{1}{2}, 1\right)$，使 $f(\eta) = \eta$;
        \item 对任意实数 $\lambda$，必存在 $\xi \in (0, \eta)$，使得 $f'(\xi) - \lambda[f(\xi) - \xi] = 1$.
        \item 存在一个 $\xi \in (0, \eta)$，使得 $f'(\xi) = f(\xi) - \xi + 1$.
    \end{enumerate}
}

\oneproblem{
    设函数 $f(x), g(x)$ 在 $[a,b]$ 上连续，在 $(a,b)$ 上内二阶可导且存在相等的最大值，又 $f(a) = g(a)$, $f(b) = g(b)$，证明:
    \begin{enumerate}[label=(\arabic*)]
        \item 存在 $\eta \in (a,b)$，使得 $f(\eta) = g(\eta)$;
        \item 存在 $\xi \in (a,b)$，使得 $f''(\xi) = g''(\xi)$.
    \end{enumerate}
}

\oneproblem{
    设函数 $f(x)$ 在区间 $[0,1]$ 上具有二阶导数，且 $f(1) > 0$, $\lim\limits_{x \to 0^+} \dfrac{f(x)}{x} < 0$. 证明:
    \begin{enumerate}[label=(\arabic*)]
        \item 方程 $f(x) = 0$ 在区间 $(0,1)$ 内至少存在一个实根;
        \item 方程 $f(x)f''(x) + [f'(x)]^2 = 0$ 在区间 $(0,1)$ 内至少存在两个不同实根.
    \end{enumerate}
}

\oneproblem{
    设函数 $f(x)$ 在 $[a,b]$ 上连续，且存在 $c \in (a,b)$ 使得 $f'(c) = 0$，证明: $\exists \xi \in (a,b)$ 使得 $f'(\xi) = \dfrac{f(\xi) - f(a)}{b - a}$.
}

\oneproblem{
    设 $f(x)$ 在 $[0,1]$ 上连续，在 $(0,1)$ 内可导，且 $f(1) = 0$，证明至少存在一点 $\xi \in (0,1)$，使 $f'(\xi) = -\dfrac{2f(\xi)}{\xi}$.
}

\oneproblem{
    设 $f(x)$ 在 $(a,+\infty)$ 内有二阶导数，且 $f(a+1) = 0$, $\lim\limits_{x \to a^+} f(x) = 0$, $\lim\limits_{x \to +\infty} f(x) = 0$. 求证: 存在 $\xi \in (a,+\infty)$ 使 $f''(\xi) = 0$.
}

\oneproblem{
    设 $f(x)$ 在 $[a,b]$ 上具有二阶导数，且 $f(a) = f(b) = 0$, $f'(a)f'(b) > 0$. 证明: 存在点 $\xi \in (a,b)$，使得 $f''(\xi) = 0$.
}

\oneproblem{
    设函数 $f(x)$ 在 $[a,b]$ 上连续，在 $(a,b)$ 内二阶可导，且 $f(a) = f(b) = 0$, $\displaystyle \int_a^b f(x)\mathrm{d}x = 0$.
    \begin{enumerate}[label=(\arabic*)]
        \item 证明: 存在互不相同的点 $x_1, x_2 \in (a,b)$，使得 $f'(x_i) = f(x_i), \ i=1,2$;
        \item 证明: 存在 $\xi \in (a,b)$, $\xi \ne x_i, \ i=1,2$，使得 $f''(\xi) = f(\xi)$.
    \end{enumerate}
}

% ======================================================
% 第 14 天：拉格朗日中值定理及其应用
% ======================================================
\setday{14}{拉格朗日中值定理及其应用}

\oneproblem{
    若函数 $f(x)$ 在区间 $(a,b)$ 内可导，$x_1, x_2$ 是区间内任意两点，且 $x_1 < x_2$，则至少存在一点 $\xi$，使得 ( )
    \begin{enumerate}[label=(\Alph*)]
        \item $f(b) - f(a) = f'(\xi)(b-a), a < \xi < b$
        \item $f(b) - f(x_1) = f'(\xi)(b-x_1), x_1 < \xi < b$
        \item $f(x_2) - f(x_1) = f'(\xi)(x_2-x_1), x_1 < \xi < x_2$
        \item $f(x_2) - f(a) = f'(\xi)(x_2-a), a < \xi < x_2$
    \end{enumerate}
}

\oneproblem{
    设 $f(x)$ 处处可导，则( )
    \begin{enumerate}[label=(\Alph*)]
        \item 当 $\lim\limits_{x \to -\infty} f(x) = -\infty$ 时，必有 $\lim\limits_{x \to -\infty} f'(x) = -\infty$
        \item 当 $\lim\limits_{x \to -\infty} f'(x) = -\infty$ 时，必有 $\lim\limits_{x \to -\infty} f(x) = -\infty$
        \item 当 $\lim\limits_{x \to +\infty} f(x) = +\infty$ 时，必有 $\lim\limits_{x \to +\infty} f'(x) = +\infty$
        \item 当 $\lim\limits_{x \to +\infty} f'(x) = +\infty$ 时，必有 $\lim\limits_{x \to +\infty} f(x) = +\infty$
    \end{enumerate}
}

\oneproblem{
    设函数 $y = f(x)$ 在 $(0,+\infty)$ 内有界且可导，则( )
    \begin{enumerate}[label=(\Alph*)]
        \item 当 $\lim\limits_{x \to +\infty} f(x) = 0$ 时，必有 $\lim\limits_{x \to +\infty} f'(x) = 0$
        \item 当 $\lim\limits_{x \to +\infty} f'(x)$ 存在时，必有 $\lim\limits_{x \to +\infty} f'(x) = 0$
        \item 当 $\lim\limits_{x \to 0^+} f(x) = 0$ 时，必有 $\lim\limits_{x \to 0^+} f'(x) = 0$
        \item 当 $\lim\limits_{x \to 0^+} f'(x)$ 存在时，必有 $\lim\limits_{x \to 0^+} f'(x) = 0$
    \end{enumerate}
}

\oneproblem{
    以下四个命题中，正确的是( )
    \begin{enumerate}[label=(\Alph*)]
        \item 若 $f'(x)$ 在 $(0,1)$ 内连续，则 $f(x)$ 在 $(0,1)$ 内有界
        \item 若 $f(x)$ 在 $(0,1)$ 内连续，则 $f(x)$ 在 $(0,1)$ 内有界
        \item 若 $f'(x)$ 在 $(0,1)$ 内有界，则 $f(x)$ 在 $(0,1)$ 内有界
        \item 若 $f(x)$ 在 $(0,1)$ 内有界，则 $f'(x)$ 在 $(0,1)$ 内有界
    \end{enumerate}
}

\oneproblem{
    求极限 $\lim\limits_{x \to \infty} x\left[\sin\ln\left(1+\dfrac{3}{x}\right) - \sin\ln\left(1+\dfrac{1}{x}\right)\right]$.
}

\oneproblem{
    设 $f(x)$ 在开区间 $(a,b)$ 内可导，且 $f'(x) > 0$，则函数 $f(x)$ 在开区间 $(a,b)$ 内严格单调增加.
}

\oneproblem{
    设不恒为常数的函数 $f(x)$ 在闭区间 $[a,b]$ 上连续，在开区间 $(a,b)$ 内可导，且 $f(a) = f(b)$. 证明: 在 $(a,b)$ 内至少存在一点 $\xi$，使 $f'(\xi) > 0$.
}

\oneproblem{
    求证: 当 $x \ge 1$ 时，$\arctan x - \dfrac{1}{2}\arccos\dfrac{2x}{1+x^2} = \dfrac{\pi}{4}$.
}

\oneproblem{
    设 $f(x)$ 在区间 $[a,b]$ 上连续，在 $(a,b)$ 内可导，证明: 在 $(a,b)$ 内至少存在一点 $\xi$，使 $\dfrac{bf(b) - af(a)}{b - a} = f(\xi) + \xi f'(\xi)$.
}

\oneproblem{
    设函数 $f(x)$ 在 $[a,b]$ 上连续，在 $(a,b)$ 内可导，且 $f(a) = f(b) = 1$，试证存在 $\xi, \eta \in (a,b)$，使得 $e^{\eta - \xi}[f(\eta) + f'(\eta)] = 1$.
}

\oneproblem{
    已知 $f(x)$ 在 $(-\infty,+\infty)$ 内可导，且 $\lim\limits_{x \to \infty} f'(x) = e$, $\lim\limits_{x \to \infty} \left(\dfrac{x+c}{x-c}\right)^x = \lim\limits_{x \to \infty} [f(x) - f(x-1)]$. 求 $c$ 的值.
}

\oneproblem{
    已知函数 $f(x)$ 在 $[0,1]$ 上连续，在 $(0,1)$ 内可导，且 $f(0) = 0, f(1) = 1$. 证明:
    \begin{enumerate}[label=(\arabic*)]
        \item 存在 $\xi \in (0,1)$, 使得 $f(\xi) = 1 - \xi$.
        \item 存在两个不同的点 $\eta, \zeta \in (0,1)$，使得 $f'(\eta)f'(\zeta) = 1$.
    \end{enumerate}
}

\oneproblem{
    设函数 $f(x)$ 在闭区间 $[0,1]$ 上连续，在开区间 $(0,1)$ 内可导，且 $f(0) = 0, f(1) = \dfrac{1}{3}$.\\证明: 存在 $\xi \in \left(0,\dfrac{1}{2}\right), \eta \in \left(\dfrac{1}{2},1\right)$，使得 $f'(\xi) + f'(\eta) = \xi^2 + \eta^2$.
}

\oneproblem{
    设函数 $f(x)$ 在 $[0,2]$ 上具有连续导数，$f(0) = f(2) = 0, M = \max\limits_{x \in [0,2]} \{|f(x)|\}$. 证明:
    \begin{enumerate}[label=(\arabic*)]
        \item 存在 $\xi \in (0,2)$，使得 $|f'(\xi)| \ge M$;
        \item 若对任意的 $x \in (0,2), |f'(x)| \le M$，则 $M = 0$.
    \end{enumerate}
}

\oneproblem{
    设 $f(x)$ 在 $[0,1]$ 上可导，且 $f(0) = 0, f(1) = 1$. 证明: 对任意正数 $a,b$，必存在 $(0,1)$ 内的两个数 $\xi$ 和 $\eta$，使 $\dfrac{a}{f'(\xi)} + \dfrac{b}{f'(\eta)} = a + b$.
}

\oneproblem{
    设函数 $f(x)$ 在 $[0,1]$ 上具有二阶导数，且满足 $f(0) = 0, f(1) = 1, f\left(\dfrac{1}{2}\right) > \dfrac{1}{4}$. 证明:
    \begin{enumerate}[label=(\arabic*)]
        \item 至少存在一点 $\xi \in (0,1)$，使得 $f''(\xi) < 2$;
        \item 若对一切 $x \in (0,1)$，有 $f''(x) \ne 2$，则当 $x \in (0,1)$ 时，恒有 $f(x) > x^2$.
    \end{enumerate}
}

\oneproblem{
    设函数 $f(x)$ 在闭区间 $[-2,2]$ 上具有二阶导数，$|f(x)| \le 1$, $[f(0)]^2 + [f'(0)]^2 = 4$. 证明: 存在一点 $\xi \in (-2,2)$，使得 $f(\xi) + f''(\xi) = 0$.
}

\oneproblem{
    设数列 $\{x_n\}$ 满足: $x_1 > 0$, $x_n e^{x_{n+1}} = e^{x_n} - 1 \ (n=1,2,\dots)$. 证明: $\{x_n\}$ 收敛，并求 $\lim\limits_{n \to \infty} x_n$.
}

% ======================================================
% 第 15 天：柯西中值定理及其应用
% ======================================================
\setday{15}{柯西中值定理及其应用}

\oneproblem{
    设函数 $f(x)$ 在 $[a,b](0 < a < b)$ 上连续，在 $(a,b)$ 内可导. 分别用柯西中值定理、罗尔中值定理证明：\\在 $(a,b)$ 内存在一点 $\xi$，使得$f(b) - f(a) = \xi f'(\xi) \ln \dfrac{b}{a}.$
}

\oneproblem{
    分别用柯西中值定理、零值定理与罗尔中值定理证明：至少存在一点 $\xi \in (1,e)$，使 $\sin 1 = \cos\ln\xi$.
}

\oneproblem{
    设 $f(x)$ 在 $[a,b]$ 上连续，在 $(a,b)$ 内可导，且 $ab > 0$. 证明：存在 $\xi \in (a,b)$，使得
    $\dfrac{ab}{b-a} \begin{vmatrix} b & a \\ f(a) & f(b) \end{vmatrix} = \xi^2\left[f(\xi) + \xi f'(\xi)\right].$
}

\oneproblem{
    设 $0 < a < b$，证明 $\exists \xi \in (a,b)$，使得
    $ae^b - be^a = (a-b)(1-\xi)e^\xi. $
}

\oneproblem{
    证明：当 $x > 0$ 时，成立不等式
    $1 + x\ln\left(x + \sqrt{1+x^2}\right) > \sqrt{1+x^2}. $
}

\oneproblem{
    求极限 $\lim\limits_{x \to 2} \dfrac{\sin\left(x^x\right) - \sin\left(2^x\right)}{2^{x^x} - 2^{2^x}}$.
}

\oneproblem{
    设 $f(x)$ 在 $[a,b]$ 上连续， $(a,b)$ 内可导， $f'(x) \ne 0$. 
    \\证明：$\exists \xi, \eta \in (a,b)$，使得
    $\dfrac{f'(\xi)}{f'(\eta)} = \dfrac{e^b - e^a}{b-a} \cdot e^{-\eta}. $
}

\oneproblem{
    设函数 $f(x)$ 在 $[a,b]$ 上连续， $(a,b)$ 内可导 $(0 \le a < b)$，\\试证在 $(a,b)$ 内存在 $\xi, \eta$，使得
    $f'(\xi) = \dfrac{a+b}{2\eta} f'(\eta). $
}

\oneproblem{
    设函数 $f(x)$ 在 $[1,2]$ 上连续，在 $(1,2)$ 内可微，且 $f'(x) \ne 0$，
    \\证明存在 $\xi, \eta, \zeta \in (1,2)$，使得
    $\dfrac{f'(\zeta)}{f'(\xi)} = \dfrac{\xi}{\eta}.$
}

\oneproblem{
    设 $0 \le a < b \le \dfrac{\pi}{2}$，$f(x)$ 在 $[a,b]$ 上连续，在 $(a,b)$ 内可导.\\ 证明：$\exists \xi, \eta \in (a,b)$，使得
    $f'(\eta)\tan\dfrac{a+b}{2} = f'(\xi)\dfrac{\sin\eta}{\cos\xi}. $
}

\oneproblem{
    设函数 $f(x), g(x)$ 在 $[a,b]$ 上连续，且 $g(a) = g(b) = 1$，在 $(a,b)$ 内 $f(x), g(x)$ 可导，且 $g(x)+g'(x) \ne 0$. $f'(x) \ne 0$.\\ 证明：$\exists \xi, \eta \in (a,b)$，使得
    $\dfrac{f'(\xi)}{f'(\eta)} = \dfrac{e^\xi\left[g(\xi) + g'(\xi)\right]}{e^\eta}. $
}

% ======================================================
% 第 16 天：洛必达法则与未定式的极限
% ======================================================
\setday{16}{洛必达法则与未定式的极限}

\oneproblem{
    设 $f(x) = \ln^{10}x, g(x) = x, h(x) = e^{\frac{x}{10}}$，则当 $x$ 充分大时有 ( )
    \begin{multicols}{4}
    \begin{enumerate}[label=(\Alph*)]
        \item $g < h < f$ \item $h < g < f$ \item $f < g < h$ \item $g < f < h$
    \end{enumerate}
    \end{multicols}
}

\oneproblem{
    设 $\lim\limits_{x \to 0} \dfrac{\ln(1+x) - (ax + bx^2)}{x^2} = 2$，则( )
    \begin{multicols}{2}
    \begin{enumerate}[label=(\Alph*)]
        \item $a=1, b=-\dfrac{5}{2}$ \item $a=0, b=-2$
        \item $a=0, b=-\dfrac{5}{2}$ \item $a=1, b=-2$
    \end{enumerate}
    \end{multicols}
}

\oneproblem{
    当 $x \to 0$ 时, $f(x) = x - \sin ax$ 与 $g(x) = x^2\ln(1-bx)$ 是等价无穷小量，则( )
    \begin{multicols}{2}
    \begin{enumerate}[label=(\Alph*)]
        \item $a=1, b=-\dfrac{1}{6}$ \item $a=1, b=\dfrac{1}{6}$
        \item $a=-1, b=-\dfrac{1}{6}$ \item $a=-1, b=\dfrac{1}{6}$
    \end{enumerate}
    \end{multicols}
}

\oneproblem{
    设函数 $f(x) = \arctan x$，若 $f(x) = x f'(\xi)$，则 $\lim\limits_{x \to 0} \dfrac{\xi^2}{x^2} =$( )
    \begin{multicols}{4}
    \begin{enumerate}[label=(\Alph*)]
        \item $1$ \item $\dfrac{2}{3}$ \item $\dfrac{1}{2}$ \item $\dfrac{1}{3}$
    \end{enumerate}
    \end{multicols}
}

\oneproblem{
    求下列极限:
    \begin{enumerate}[label=(\arabic*), itemsep=6cm]
        \item $\lim\limits_{x \to +\infty}\dfrac{\ln\left(1+\frac{1}{x}\right)}{\arccot x}$.
        \item $\lim\limits_{x \to +\infty} \left(x + \sqrt{1+x^2}\right)^{\frac{1}{x}}$.
        \item $\lim\limits_{x \to 0} \left(\dfrac{e^x + e^{2x} + \dots + e^{nx}}{n}\right)^{\frac{1}{x}}$.
        \item $\lim\limits_{x \to \infty} \left[x - x^2\ln\left(1+\frac{1}{x}\right)\right]$.
        \item $\lim\limits_{n \to \infty} \left(n\tan\frac{1}{n}\right)^{n^2}$.
        \item $\lim\limits_{x \to +\infty} x^2\left[\arctan(x+1) - \arctan x\right]$.
        \item $\lim\limits_{n \to \infty} \left(\frac{a^{1/n} + b^{1/n} + c^{1/n}}{3}\right)^n$.
        \item $\lim\limits_{x \to 0} \left[\frac{a}{x} - \left(\frac{1}{x^2} - a^2\right)\ln(1+ax)\right](a \ne 0)$.
        \item $\lim\limits_{x \to 0} \dfrac{[\sin x - \sin(\sin x)]\sin x}{x^4}$.
    \end{enumerate}
}

\oneproblem{
    设 $f(x,y) = \dfrac{y}{1+xy} - \dfrac{1-y\sin\frac{\pi x}{y}}{\arctan x}$，求:
    \begin{enumerate}[label=(\arabic*)]
        \item $g(x) = \lim\limits_{y \to +\infty} f(x,y)$;
        \item $\lim\limits_{x \to 0^+} g(x)$.
    \end{enumerate}
}

\oneproblem{
    设 $f(x) = \begin{cases} \dfrac{g(x) - e^{-x}}{x}, & x \ne 0 \\ 0, & x=0 \end{cases}$，其中 $g(x)$ 有二阶连续导数，$g(0)=1, g'(0)=-1$.
    \begin{enumerate}[label=(\arabic*)]
        \item 求 $f'(x)$;
        \item 讨论 $f'(x)$ 在 $(-\infty,+\infty)$ 上的连续性.
    \end{enumerate}
}

\oneproblem{
    设函数 $f(x)$ 在 $x=0$ 的某邻域内有二阶连续导数，且 $f(0), f'(0), f''(0)$ 均不为零. 证明：存在唯一一组实数 $k_1,k_2,k_3$，使得
    $ \lim\limits_{h \to 0} \dfrac{k_1f(h) + k_2f(2h) + k_3f(3h) - f(0)}{h^2} = 0. $
}

% ======================================================
% 第 17 天：泰勒公式之计算篇
% ======================================================
\setday{17}{泰勒公式之计算篇}

\oneproblem{
    设当 $x \to 0$ 时，$e^x - \left(ax^2 + bx + 1\right)$ 是比 $x^2$ 高阶的无穷小，则( )
    \begin{multicols}{4}
    \begin{enumerate}[label=(\Alph*)]
        \item $a=\frac{1}{2}, b=1$ \item $a=1, b=1$ \item $a=-\frac{1}{2}, b=1$ \item $a=-1, b=1$
    \end{enumerate}
    \end{multicols}
}

\oneproblem{
    若 $\lim\limits_{x \to 0} \dfrac{\sin 6x + x f(x)}{x^3} = 0$，则 $\lim\limits_{x \to 0} \dfrac{6 + f(x)}{x^2}$ 为( )
    \begin{multicols}{4}
    \begin{enumerate}[label=(\Alph*)]
        \item $0$ \item $6$ \item $36$ \item $\infty$
    \end{enumerate}
    \end{multicols}
}

\oneproblem{
    设 $f(x) = x + a\ln(1+x) + bx\sin x, g(x) = kx^3$，若 $f(x) \sim g(x) \ (x \to 0)$，求 $a,b,k$ 值.
}

\oneproblem{
    当 $x \to 0$ 时, $1-\cos x\cos2x\cos3x$ 与 $ax^n$ 是等价无穷小量，求常数 $n$ 与 $a$ 的值.
}

\oneproblem{
    求函数 $f(x) = \dfrac{1}{2+x}$ 在 $x_0=1$ 处的 $7$ 阶带佩亚诺余项的泰勒公式，并求 $f^{(7)}(1)$.
}

\oneproblem{
    求 $f(x) = x^2\ln(1+x)$ 在 $x=0$ 处的 $n$ 阶导数 $f^{(n)}(0)(n \ge 3)$.
}

\oneproblem{
    设 $f(x) = e^x\sin 2x$，求 $f^{(4)}(0)$.
}

\oneproblem{
    求下列极限:
    \begin{enumerate}[label=(\arabic*), itemsep=6cm]
        \item $\lim\limits_{x \to 0} \dfrac{\sqrt{1+2\sin x} - x - 1}{x\ln(1+x)}$.
        \item $\lim\limits_{x \to 0} (\cos 2x + 2x\sin x)^{\frac{1}{x^4}}$.
        \item $\lim\limits_{x \to \infty} e^{-x}\left(1+\frac{1}{x}\right)^{x^2}$.
        \item $\lim\limits_{x \to +\infty} \left(\sqrt[6]{x^6+x^5} - \sqrt[6]{x^6-x^5}\right)$.
        \item $\lim\limits_{x \to \infty} \left[\left(x^3 - x^2 + \frac{x}{2}\right)e^{\frac{1}{x}} - \sqrt{x^6+1}\right]$.
        \item $\lim\limits_{n \to \infty} [n\sin(\pi n!e)]$.
    \end{enumerate}
}

\oneproblem{
    设 $f(x) = (1+x)^{\frac{1}{x}} \ (x>0)$，求证:
    $f(x) = e + Ax + Bx^2 + o(x^2) $
    并求 $A,B$.
}

\oneproblem{
    设 $f(x)$ 具有二阶连续导数，且 $f(0)=f'(0)=0, f''(0)=6$，求 $\lim\limits_{x \to 0} \dfrac{f\left(\sin^2 x\right)}{x^4}$.
}

\oneproblem{
    若 $f(x)$ 在点 $x=a$ 处可导，且 $f(a) \ne 0$，求 $\lim\limits_{n \to +\infty} \left[\dfrac{f(a+1/n)}{f(a)}\right]^n$.
}

\oneproblem{
    设函数 $y=f(x)$ 二阶可导，且 $f''(x)>0, f(0)=0, f'(0)=0$. 求 $\lim\limits_{x \to 0} \dfrac{x^3f(u)}{f(x)\sin^3 u}$，其中 $u$ 是切线在 $x$ 轴上的截距.
}

\oneproblem{
    求方程 $x^2\sin\dfrac{1}{x} = 2x - 501$ 的近似解，精确到 $0.001$.
}

\oneproblem{
    求一函数 $f(x)$，使其在任一有限区间上有界，且满足方程
    $f(x) - \dfrac{1}{2}f\left(\dfrac{x}{2}\right) = x - x^2$.
}

\oneproblem{
    设 $f(x)$ 三阶连续可导，满足 $f(0)=0, f'(0)=1, f''(0)=0, f'''(0)=-1$；又设 $a_{n+1}=f(a_n)$，$a_n \to 0$. 计算 $\lim\limits_{n \to \infty} na_n^2$.
}

% ======================================================
% 第 18 天：泰勒公式之证明篇
% ======================================================
\setday{18}{泰勒公式之证明篇}

\oneproblem{
    设 $f(x)$ 在 $x_0$ 邻域有 $n$ 阶连续导数，$f^{(k)}(x_0)=0\ (k=2,\dots,n-1), f^{(n)}(x_0)\ne 0$. 已知 $f(x_0+h)-f(x_0)=hf'(x_0+\theta h)$. 证明：$\lim\limits_{h\to 0}\theta=\dfrac{1}{\sqrt[n-1]{n}}$.
}

\oneproblem{
    设 $f(x)$ 在 $[-1,1]$ 上三阶连续可导，且 $f(-1)=0, f(1)=1, f'(0)=0$. 证明：在 $(-1,1)$ 内存在 $\xi$，使 $f'''(\xi)=3$.
}

\oneproblem{
    设 $f(x)$ 在 $[a,b]$ 上具有连续的二阶导数，证明：存在 $\xi\in(a,b)$，使得
    $ \dfrac{4}{(b-a)^2}\left[f(a)-2f\left(\dfrac{a+b}{2}\right)+f(b)\right]=f''(\xi). $
}

\oneproblem{
    设 $f(x)$ 在 $[0,1]$ 上具有二阶导数，$|f(x)|\le a, |f''(x)|\le b$. 对任意 $c \in (0,1)$，证明：$|f'(c)|\le 2a+\dfrac{b}{2}$.
}

\oneproblem{
    设 $x\in[0,2]$ 时，$|f(x)|\le1, |f''(x)|\le1$. 证明：对 $x\in[0,2]$，有 $|f'(x)|\le2$.
}

\oneproblem{
    设 $f(x)$ 在 $(a,+\infty)$ 内二阶可导，且 $f, f''$ 有界. 证明：$f'$ 在 $(a,+\infty)$ 内有界.
}

\oneproblem{
    设 $f(x)$ 在 $[a,b]$ 上连续，在 $(a,b)$ 二阶可导，$|f''(x)|\ge m>0$，且 $f(a)=f(b)=0$. 证明：$\max_{[a,b]}|f(x)|\ge\dfrac{m}{8}(b-a)^2$.
}

\oneproblem{
    设 $f(x)$ 在 $x_0$ 有 $n$ 阶导数，$f'(x_0)=\dots=f^{(n-1)}(x_0)=0, f^{(n)}(x_0)\ne 0$. 讨论 $n$ 为偶数和奇数时极值点的情况.
}

\oneproblem{
    设 $f(x)$ 在 $[a,b]$ 上有 $n$ 阶导数，$f^{(k)}(a)=f^{(k)}(b)=0\ (k=1,\dots,n-1)$. 证明：$\exists c\in(a,b)$ 使得
    $|f^{(n)}(c)|\ge\dfrac{2^{n-1}n!}{(b-a)^n}|f(b)-f(a)|$.
}

\oneproblem{
    设 $f(x+h)=f(x)+f'(x)h+\dfrac{1}{2}f''(x+\theta h)h^2$，其中 $\theta$ 是常数. 证明 $f$ 是不超过三次的多项式.
}

\oneproblem{
    设 $\alpha>0, f \in C^1[0,1], f^{(n)}(0)=0$. 存在 $C>0$ 使得 $|x^\alpha f'(x)|\le C|f(x)|$.
    讨论 $\alpha=1$ 和 $\alpha>1$ 时 $f(x)$ 是否恒为零.
}

\oneproblem{
    设 $f(0)=0$ 且在 $0$ 处可导. 证明：$\lim\limits_{n\to\infty}\sum_{k=1}^n f\left(\dfrac{k}{n^2}\right)=\dfrac{f'(0)}{2}$.
}

\oneproblem{
    证明对任何 $\alpha \in \mathbb{R}$，存在 $a_k \in \{-1,1\}$ 满足 $\lim\limits_{n\to+\infty}\left(\sum\limits_{k=1}^n\sqrt{n+a_k}-n^{\frac{3}{2}}\right)=\alpha$.
}

% ======================================================
% 第 19 天：函数的单调性与极值
% ======================================================
\setday{19}{函数的单调性与极值}

\oneproblem{
    设 $\lim\limits_{x\to a}\dfrac{f(x)-f(a)}{(x-a)^2}=-1$，则在 $x=a$ 处 ( )
    \begin{multicols}{4}
    \begin{enumerate}[label=(\Alph*)]
        \item $f'(a)\ne0$ \item 取得极大值 \item 取得极小值 \item 导数不存在
    \end{enumerate}
    \end{multicols}
}

\oneproblem{
    设 $g''(x)<0$，若 $g(x_0)=a$ 是极值，则 $f(g(x))$ 在 $x_0$ 取极大值的充分条件是 ( )
    \begin{multicols}{4}
    \begin{enumerate}[label=(\Alph*)]
        \item $f'(a)<0$ \item $f'(a)>0$ \item $f''(a)<0$ \item $f''(a)>0$
    \end{enumerate}
    \end{multicols}
}

\oneproblem{
    设 $f'(x)>f(x)>0$，则在 $[-2,2]$ 上 ( )
    \begin{multicols}{4}
    \begin{enumerate}[label=(\Alph*)]
        \item $\frac{f(-2)}{f(-1)}>1$ \item $\frac{f(0)}{f(-1)}>e$
        \item $\frac{f(1)}{f(-1)}<e^2$ \item $\frac{f(2)}{f(-1)}<e^3$
    \end{enumerate}
    \end{multicols}
}

\oneproblem{
    已知 $f(x)=\begin{cases}x^{2x},&x>0,\\xe^x+1,&x\le0.\end{cases}$ 求 $f'(x)$ 及 $f(x)$ 的极值.
}

\oneproblem{
    设 $f(0)=f(1)=0, f''(x)+2f'(x)+f(x)>0$. 证明：$f(x)\le0, \forall x\in(0,1)$.
}

\oneproblem{
    设 $f(x)=nx(1-x)^n$，求 $[0,1]$ 上的最大值 $M(n)$ 及 $\lim\limits_{n\to\infty}M(n)$.
}

\oneproblem{
    设 $y=y(x)$ 由 $x^3+3x^2y-2y^3=2$ 确定，求 $y(x)$ 的极值.
}

\oneproblem{
    证明：$x-\dfrac{x^2}{2}<\ln(1+x)<x-\dfrac{x^2}{2(1+x)}, x>0$.
}

\oneproblem{
    证明：$e^{-x}+\sin x<1+\dfrac{x^2}{2}, 0<x<\dfrac{\pi}{2}$.
}

\oneproblem{
    设 $n>1$，证明：$\dfrac{1}{2ne}<\dfrac{1}{e}\left(1-\dfrac{1}{n}\right)^n<\dfrac{1}{ne}$.
}

\oneproblem{
    设 $b>a>e$，证明：$a^b>b^a$.
}

\oneproblem{
    设 $0<a<b$，证明：$\dfrac{2a}{a^2+b^2}<\dfrac{\ln b-\ln a}{b-a}<\dfrac{1}{\sqrt{ab}}$.
}

\oneproblem{
    设 $0<x<\dfrac{\pi}{2}$，证明：$\dfrac{4}{\pi^2}<\dfrac{1}{x^2}-\dfrac{1}{\tan^2x}<\dfrac{2}{3}$.
}

\oneproblem{
    求使得 $e<\left(1+\frac{1}{n}\right)^{n+\beta}$ 恒成立的最小 $\beta$.
}

\oneproblem{
    证明不等式：$2^n\ge1+n\sqrt{2^{n-1}} \ (n \ge 1)$.
}

\oneproblem{
    设计圆柱体容器，容积 $V$，底材料费 $a$/面积，侧材料费 $b$/面积. 求最省钱的高径比.
}

\oneproblem{
    设 $P_A, P_B$ 为单位圆外切 $n$ 边形和切点多边形周长. 证：$P_A^{\frac{1}{3}}P_B^{\frac{2}{3}}>2\pi$.
}

\oneproblem{
    设 $\alpha>1$. 证不存在正可导函数满足 $f'(x)\ge f^\alpha(x), x\in[0,+\infty)$.
}

% ======================================================
% 第 20 天：函数曲线的凹凸性与拐点
% ======================================================
\setday{20}{函数曲线的凹凸性与拐点}

\oneproblem{
    设 $f''(x)+[f'(x)]^2=x, f'(0)=0$，则( )
    \begin{multicols}{4}
    \begin{enumerate}[label=(\Alph*)]
        \item $f(0)$ 是极大值 \item $f(0)$ 是极小值 \item $(0,f(0))$ 是拐点 \item 均不是
    \end{enumerate}
    \end{multicols}
}

\oneproblem{
    曲线 $y=(x-1)^2(x-3)^2$ 的拐点个数为( )
    \begin{multicols}{4}
    \begin{enumerate}[label=(\Alph*)]
        \item 0 \item 1 \item 2 \item 3
    \end{enumerate}
    \end{multicols}
}

\oneproblem{
    设 $f'(x_0)=f''(x_0)=0, f'''(x_0)>0$，则 ( )
    \begin{multicols}{4}
    \begin{enumerate}[label=(\Alph*)]
        \item $f'(x_0)$ 是极大值 \item $f(x_0)$ 是极大值 \item $f(x_0)$ 是极小值 \item $(x_0,f(x_0))$ 是拐点
    \end{enumerate}
    \end{multicols}
}

\oneproblem{
    确定 $y=k(x^2-3)^2$ 中 $k$ 的值，使拐点处法线过原点.
}

\oneproblem{
    判断曲线 $y\ln y-x+y=0$ 在点 $(1,1)$ 附近的凹凸性.
}

\oneproblem{
    参数方程 $x=\frac{1}{3}t^3+t+\frac{1}{3}, y=\frac{1}{3}t^3-t+\frac{1}{3}$，求极值、凹凸性及拐点.
}

\oneproblem{
    $f(a)=0, f', f''>0$. 证：切线与 $x$ 轴交点 $x_0$ 满足 $a<x_0<b$.
}

\oneproblem{
    证明：$e^{\frac{a+b}{2}}\le\dfrac{1}{2}(e^a+e^b)$.
}

\oneproblem{
    证明琴生不等式：若 $f$ 为凸函数，则 $f\left(\sum \lambda_i x_i\right)\le\sum \lambda_i f(x_i)$.
}

\oneproblem{
    设 $f(0)=0, f(1)=1, f''(x)<0$. 证明：$f(x)\ge x, \forall x\in[0,1]$.
}

\oneproblem{
    已知 $f>0, f''f-(f')^2\ge 0$. 
    \begin{enumerate}[label=(\arabic*)]
        \item 证明: $f(x_1)f(x_2)\ge f^2\left(\frac{x_1+x_2}{2}\right)$;
        \item 若 $f(0)=1$，证明 $f(x)\ge e^{f'(0)x}$.
    \end{enumerate}
}